\chapter{Transforming Functions}

Knowing the catalogue of standard shapes is worth much more once it is
combined with the second half of the skill: recognising that an unfamiliar
formula is a familiar function in disguise. The graph of
$y=-2\mathrm{e}^{-(x-3)}+5$ never has to be plotted point by point. It is
the exponential, mirrored twice, stretched, and moved, and all of this is
read off the formula without any arithmetic.

Three related skills belong here. \textbf{Transforming}: what the constants
in $a\,f\bigl(b(x-c)\bigr)+d$ do to a graph, with the standing trap that
everything \emph{inside} the function acts in the opposite direction and by
the reciprocal factor. \textbf{Composing and decomposing}: building a
function out of simpler ones, and taking one apart again, which is exactly
the reading that makes the chain rule mechanical later. And
\textbf{inverting}: swapping the roles of $x$ and $y$, the operation that
turns an exponential into a logarithm and a demand curve into an inverse
demand curve.

The chapter ends with the reverse question, the one an economist meets
most often: not ``what does this formula look like'', but ``which formula has
the shape I need?'' Given a vertex, an asymptote, a saturation level or two
data points, determine the parameters.

\section{Core Rules}

\begin{rulebox}{The four elementary moves}
Start from a known graph $y=f(x)$ and read the general form
\[
  y=a\,f\bigl(b(x-c)\bigr)+d .
\]
\begin{itemize}[leftmargin=1.6em,itemsep=0.3ex]
  \item \textbf{$d$: vertical shift.} $f(x)+d$ moves the graph \emph{up}
        by $d$ (down if $d<0$). Outside the function, and in the direction
        you expect.
  \item \textbf{$c$: horizontal shift.} $f(x-c)$ moves the graph
        \emph{right} by $c$. Inside the function, and \emph{against}
        the sign: $f(x+2)$ moves two units to the \emph{left}.
  \item \textbf{$a$: vertical scaling.} $a\,f(x)$ stretches the graph
        vertically by the factor $|a|$ ($0<|a|<1$ compresses), and
        \textbf{reflects it in the $x$-axis} if $a<0$. Roots stay put;
        every $y$-value is multiplied by $a$.
  \item \textbf{$b$: horizontal scaling.} $f(bx)$ \emph{compresses} the
        graph horizontally by the factor $|b|$ --- again the reciprocal of
        what the number suggests --- and \textbf{reflects it in the
        $y$-axis} if $b<0$. The $y$-intercept stays put.
\end{itemize}
Order matters: inside operations act on $x$ before $f$ is applied, outside
operations act on the result. Apply them as
\emph{horizontal scaling and shifting first, then $f$, then vertical
scaling and shifting}. And factor before reading: $f(2x-6)=f\bigl(2(x-3)\bigr)$
is a compression by $2$ together with a shift of $3$ to the right, not a
shift of $6$.
\end{rulebox}

\begin{gallery}
\begin{tikzpicture}
  \begin{axis}[fnaxis, xmin=-2.6, xmax=4.6, ymin=-1.2, ymax=4.6,
      xtick={-2,2,4}, ytick={2,4}, title={$y=(x-2)^{2}+1$}]
    \addplot[gray!60, dashed, thin, samples=101, domain=-2.1:2.1]{x^2};
    \addplot[domain=0.15:3.85]{(x-2)^2+1};
  \end{axis}
\end{tikzpicture}\hfill
\begin{tikzpicture}
  \begin{axis}[fnaxis, xmin=-2.6, xmax=2.6, ymin=-1.2, ymax=4.6,
      xtick={-2,2}, ytick={2,4}, title={$y=2x^{2}$}]
    \addplot[gray!60, dashed, thin, samples=101, domain=-2.1:2.1]{x^2};
    \addplot[domain=-1.5:1.5]{2*x^2};
  \end{axis}
\end{tikzpicture}\hfill
\begin{tikzpicture}
  \begin{axis}[fnaxis, xmin=-3.4, xmax=3.4, ymin=-1.2, ymax=4.6,
      xtick={-2,2}, ytick={2,4}, title={$y=\mathrm{e}^{-x}$}]
    \addplot[gray!60, dashed, thin, domain=-3.2:1.5]{exp(x)};
    \addplot[domain=-1.5:3.2]{exp(-x)};
  \end{axis}
\end{tikzpicture}
\gcap{Shift right and up \quad\textbullet\quad vertical stretch
\quad\textbullet\quad reflection in the $y$-axis. The original is dashed.}
\end{gallery}

\begin{rulebox}{Rewrite first, then read: other bases, negative inner factors, higher powers}
The four moves can be read off only once the formula has the shape
$a\,f\bigl(b(x-c)\bigr)+d$ for a base function you know. Four rewritings
get it there.
\begin{itemize}[leftmargin=1.6em,itemsep=0.3ex]
  \item \textbf{Exponentials with any base.} For $q>0$, $q\neq1$, the graph
        of $q^{x}$ passes through $(0,1)$ and has the asymptote $y=0$; it
        rises if $q>1$ and falls if $0<q<1$. Since $q^{-x}=(1/q)^{x}$, a
        negative exponent is a reflection in the $y$-axis, turning a rising
        exponential into a falling one. Collect everything into a single
        base before sketching:
        \[
          \frac{a^{x}}{b^{x}}=\Bigl(\frac{a}{b}\Bigr)^{x},\qquad
          a^{2x}=\bigl(a^{2}\bigr)^{x},\qquad
          \text{e.g. }\ \frac{3^{2x}}{2^{x}}=\Bigl(\frac{9}{2}\Bigr)^{x}.
        \]
        Then $y=c\,q^{x}+d$ has domain $\R$, the horizontal asymptote
        $y=d$, and range $(d,\infty)$ if $c>0$, $(-\infty,d)$ if $c<0$.
  \item \textbf{A negative factor inside.} In $\sqrt{k-mx}$ with $m>0$,
        factor out $-m$, not $m$:
        \[
          \sqrt{k-mx}=\sqrt{-m\Bigl(x-\frac{k}{m}\Bigr)} .
        \]
        The factor $-m$ is a reflection in the $y$-axis together with a
        horizontal compression by $m$; then comes a shift \emph{right} by
        $k/m$. The domain is $x\le k/m$: the curve starts at $x=k/m$ and
        runs to the \emph{left}. The same factoring handles $\ln(k-mx)$,
        with domain $x<k/m$ and the asymptote $x=k/m$.
  \item \textbf{Logarithms with an inner factor.}
        $c\ln(ax+b)+d=c\ln\bigl(a(x+\tfrac{b}{a})\bigr)+d$: domain
        $ax+b>0$, vertical asymptote $x=-b/a$, range $\R$. For $a>0$ there
        is a second reading, $c\ln a+c\ln\bigl(x+\tfrac{b}{a}\bigr)+d$: the
        horizontal compression by $a$ is the same as a vertical shift by
        $c\ln a$.
  \item \textbf{Higher powers in a denominator.} $y=\dfrac{c}{(x-a)^{n}}+b$
        has the asymptotes $x=a$ and $y=b$. For \emph{even} $n$ the
        denominator is positive on both sides of the pole, so both branches
        lie on the \emph{same} side of $y=b$ (above it if $c>0$, below if
        $c<0$), and the range is $(b,\infty)$ or $(-\infty,b)$. For
        \emph{odd} $n$ the denominator changes sign at $x=a$, the branches
        lie on \emph{opposite} sides, as for $1/x$, and the range is
        $\R\setminus\{b\}$.
\end{itemize}
Factor powers too: $(2x-6)^{2}=4(x-3)^{2}$, so
$\dfrac{8}{(2x-6)^{2}}=\dfrac{2}{(x-3)^{2}}$, with the pole at $x=3$.
\end{rulebox}

\begin{gallery}
\begin{tikzpicture}
  \begin{axis}[fnaxis, xmin=-3.4, xmax=3.4, ymin=-1.2, ymax=4.6,
      xtick={-2,2}, ytick={2,4}, title={$y=2^{-x}=\bigl(\tfrac12\bigr)^{x}$}]
    \addplot[gray!60, dashed, thin, domain=-3.2:2.1]{2^x};
    \addplot[domain=-2.1:3.2]{2^(-x)};
  \end{axis}
\end{tikzpicture}\hfill
\begin{tikzpicture}
  \begin{axis}[fnaxis, xmin=-3.4, xmax=4.4, ymin=-1.2, ymax=4.6,
      xtick={-2,2,4}, ytick={2,4}, title={$y=\sqrt{4-2x}$}]
    \addplot[gray!60, dashed, thin, samples=101, domain=0:4.2]{sqrt(x)};
    \addplot[samples=201, domain=-3.2:2]{sqrt(4-2*x)};
  \end{axis}
\end{tikzpicture}\hfill
\begin{tikzpicture}
  \begin{axis}[fnaxis, xmin=-2.4, xmax=3.4, ymin=-3.2, ymax=4.6,
      xtick={-2,2}, ytick={-2,2,4},
      title={$\tfrac{1}{(x-1)^{2}}$ and $\tfrac{1}{(x-1)^{3}}$}]
    \addplot[samples=101, domain=-2.3:0.53]{1/(x-1)^2};
    \addplot[samples=101, domain=1.47:3.3]{1/(x-1)^2};
    \addplot[gray!60, dashed, thin, samples=101, domain=-2.3:0.31]{1/(x-1)^3};
    \addplot[gray!60, dashed, thin, samples=101, domain=1.62:3.3]{1/(x-1)^3};
  \end{axis}
\end{tikzpicture}
\gcap{A negative exponent reflects in the $y$-axis \quad\textbullet\quad
$\sqrt{-2(x-2)}$ starts at $x=2$ and runs left \quad\textbullet\quad even
power (solid): both branches above; odd power (dashed): opposite sides.}
\end{gallery}

\begin{rulebox}{Absolute values: two different operations}
\begin{itemize}[leftmargin=1.6em,itemsep=0.3ex]
  \item $y=|f(x)|$ acts on the \textbf{output}: everything below the
        $x$-axis is folded \emph{up}, everything above stays. Roots become
        points where the graph touches the axis and turns.
  \item $y=f(|x|)$ acts on the \textbf{input}: the part of the graph with
        $x\geq0$ is kept and \emph{mirrored} to the left, and whatever the
        function did for $x<0$ is discarded. The result is always even.
\end{itemize}
The two coincide only by accident. Note also $|ax+b|=|a|\,
\bigl|x+\tfrac{b}{a}\bigr|$, which turns any such expression into a shifted
and stretched V.
\end{rulebox}

\begin{gallery}
\begin{tikzpicture}
  \begin{axis}[fnaxis, xmin=-2.6, xmax=2.6, ymin=-2.6, ymax=2.6,
      xtick={-2,2}, ytick={-2,2}, title={$y=f(x)=x^{3}-3x$}]
    \addplot[domain=-2.1:2.1]{x^3-3*x};
  \end{axis}
\end{tikzpicture}\hfill
\begin{tikzpicture}
  \begin{axis}[fnaxis, xmin=-2.6, xmax=2.6, ymin=-2.6, ymax=2.6,
      xtick={-2,2}, ytick={2}, title={$y=|f(x)|$}]
    \addplot[domain=-2.1:2.1]{abs(x^3-3*x)};
  \end{axis}
\end{tikzpicture}\hfill
\begin{tikzpicture}
  \begin{axis}[fnaxis, xmin=-2.6, xmax=2.6, ymin=-2.6, ymax=2.6,
      xtick={-2,2}, ytick={-2,2}, title={$y=f(|x|)$}]
    \addplot[domain=-2.1:2.1]{abs(x)^3-3*abs(x)};
  \end{axis}
\end{tikzpicture}
\end{gallery}

\begin{rulebox}{Composing and decomposing}
$(f\circ g)(x)=f\bigl(g(x)\bigr)$: the \textbf{inner} function $g$ is
applied first. To \emph{decompose} a formula, ask what you would compute
first with a calculator; that is the inner function:
\[
  \begin{aligned}
    \sqrt{3x+1}\;&:\quad g(x)=3x+1, \quad f(u)=\sqrt{u},\\
    \mathrm{e}^{x^{2}}\;&:\quad g(x)=x^{2}, \quad f(u)=\mathrm{e}^{u},\\
    (\ln x)^{3}\;&:\quad g(x)=\ln x, \quad f(u)=u^{3}.
  \end{aligned}
\]
Decompositions are not unique, but the ``innermost bracket first'' reading
is the one the chain rule wants. Note that all four elementary moves are
themselves compositions with linear functions: $f(b(x-c))$ is $f$ composed
with a linear inner function.
\end{rulebox}

\begin{rulebox}{Inverse functions}
$f^{-1}$ undoes $f$: \;$f^{-1}\bigl(f(x)\bigr)=x$ and
$f\bigl(f^{-1}(y)\bigr)=y$. It exists exactly when $f$ is
\textbf{injective}; if it is not, restrict the domain first (as with
$x^{2}$ on $[0,\infty)$). Three consequences and one recipe:
\begin{itemize}[leftmargin=1.6em,itemsep=0.3ex]
  \item $D(f^{-1})=W(f)$ and $W(f^{-1})=D(f)$: domain and range swap.
  \item The graph of $f^{-1}$ is the graph of $f$ \textbf{mirrored in the
        line $y=x$}; the point $(p,q)$ becomes $(q,p)$.
  \item $f^{-1}$ is \emph{not} $1/f$.
  \item \textbf{Recipe:} write $y=f(x)$, solve for $x$, then rename the
        variables. \;$y=\mathrm{e}^{2x}\Rightarrow \ln y=2x\Rightarrow
        x=\tfrac12\ln y$, so $f^{-1}(x)=\tfrac12\ln x$ with
        $D(f^{-1})=(0,\infty)$.
\end{itemize}
This is the operation behind every \emph{inverse demand} function: the
demand $q(p)$ solved for $p$ gives $p(q)$, the same curve read along the
other axis.
\end{rulebox}

\begin{gallery}
\begin{tikzpicture}
  \begin{axis}[fnaxis, xmin=-3.2, xmax=4.4, ymin=-3.2, ymax=4.4,
      xtick={-2,2,4}, ytick={-2,2,4},
      title={$\mathrm{e}^{x}$ and $\ln x$}]
    \addplot[domain=-3:1.45]{exp(x)};
    \addplot[domain=0.05:4.2]{ln(x)};
    \addplot[gray!60, dashed, thin, samples=2, domain=-3:4.2]{x};
  \end{axis}
\end{tikzpicture}\hfill
\begin{tikzpicture}
  \begin{axis}[fnaxis, xmin=-0.8, xmax=4.4, ymin=-0.8, ymax=4.4,
      xtick={2,4}, ytick={2,4},
      title={$x^{2}$ on $[0,\infty)$ and $\sqrt{x}$}]
    \addplot[domain=0:2.05]{x^2};
    \addplot[domain=0:4.2]{sqrt(x)};
    \addplot[gray!60, dashed, thin, samples=2, domain=0:4.2]{x};
  \end{axis}
\end{tikzpicture}
\gcap{Each pair is symmetric about the dashed line $y=x$.}
\end{gallery}

\begin{rulebox}{Reverse engineering: choosing the parameters}
Given the shape you want, pick the family and fit the constants.
\begin{itemize}[leftmargin=1.6em,itemsep=0.3ex]
  \item \textbf{Prescribed vertex:} use $y=a(x-h)^{2}+k$ with the vertex
        $(h,k)$, and determine $a$ from one further point.
  \item \textbf{Prescribed asymptotes:} $y=\dfrac{k}{x-c}+d$ has the
        vertical asymptote $x=c$ and the horizontal one $y=d$; a rational
        $\dfrac{ax+b}{x-c}$ has $y=a$ as horizontal asymptote.
  \item \textbf{Prescribed saturation:} $y=S\bigl(1-\mathrm{e}^{-kx}\bigr)$
        starts at $0$ and rises to the ceiling $S$; $y=S\mathrm{e}^{-kx}$
        starts at $S$ and decays to $0$.
  \item \textbf{Two data points, exponential family:} $y=A\,q^{x}$ with
        $A=y(0)$, and $q$ from the ratio of the two values. Halving every
        $T$ units means $q=2^{-1/T}$.
  \item \textbf{Economic shifts:} a change in fixed cost shifts $C(q)$
        \emph{vertically}; a per-unit tax $t$ adds $tq$ and therefore
        changes the \emph{slope}, not the intercept; measuring the input in
        thousands instead of units rescales $x$ horizontally.
\end{itemize}
\end{rulebox}

\section{Worked Example}

\begin{worked}{Describe the moves taking $y=f(x)$ to $y=-3f(2x-4)+1$, apply them to $f(x)=\sqrt{x}$, and invert $g(x)=2\ln(x-1)+3$}
\wstep{Step 1. Factor the inside before saying anything about the shift.}
The inside is $2x-4$, and it must be written as
\[
  2x-4=2(x-2)
\]
before it can be read. In this form the horizontal instructions are
unambiguous: \textbf{compress by a factor $2$}, and \textbf{shift right by
$2$} --- not by $4$.
\wtrap{This is the chapter's central trap. $f(2x-4)$ shifts by $2$, because
what matters is the value of $x$ that makes the argument zero. Reading the
shift straight off the $-4$ is the most common error with
transformations.}

\wstep{Step 2. Separate inside from outside, and get the directions right.}
\begin{center}\footnotesize
\begin{tabular}{lll}
  \hline
  where & operation & effect on the graph\\
  \hline
  inside, $2(\cdot)$ & multiply $x$ by $2$ & horizontal compression by $2$ \\
  inside, $(\cdot-2)$ & subtract $2$ from $x$ & shift \emph{right} $2$ \\
  outside, $-3(\cdot)$ & multiply by $-3$ & vertical stretch $3$, flip in the $x$-axis\\
  outside, $+1$ & add $1$ & shift \emph{up} $1$ \\
  \hline
\end{tabular}
\end{center}
Inside changes act \textbf{opposite} to how they read (subtracting moves
right, multiplying compresses); outside changes act \textbf{as} they read.
The reason is that inside operations act on the input before $f$ sees it,
so to keep the output the same you must compensate.

\wstep{Step 3. Apply the moves in the right order.}
Inside: shift and scale, in the order that the factored form dictates ---
compress first, then shift, or equivalently shift the compressed picture by
$2$. Outside: stretch and reflect first, then translate. Applying the
vertical shift before the stretch would multiply the shift by $-3$ as well,
landing at $-3f(\cdot)-3$ instead.

\wstep{Step 4. Apply it to $f(x)=\sqrt{x}$.}
\[
  y=-3\sqrt{2x-4}+1=-3\sqrt{2(x-2)}+1 .
\]
The domain moves with the graph: $\sqrt{\cdot}$ needs $2x-4\ge0$, so
$x\ge2$. The starting point $(0,0)$ of $\sqrt{x}$ goes to $(2,1)$: $x=2$
because the shift is right by~$2$, and $y=1$ because $-3\cdot0+1=1$. From
there the curve falls (the reflection), so its range is $(-\infty,1]$.

\wstep{Step 5. Track one more point, as a check.}
The point $(4,2)$ on $y=\sqrt{x}$: solve $2x-4=4$ for the new abscissa,
giving $x=4$; the new ordinate is $-3\cdot2+1=-5$. So $(4,2)\mapsto(4,-5)$,
and substituting $x=4$ into the final formula gives
$-3\sqrt{4}+1=-5$. \checkmark

\wstep{Step 6. Now the inverse of $g(x)=2\ln(x-1)+3$: undo the operations in reverse order.}
Forward, $g$ does: subtract $1$, take $\ln$, double, add $3$. Backwards:
subtract $3$, halve, exponentiate, add $1$. In symbols, set $y=g(x)$ and
solve for $x$:
\[
  y-3=2\ln(x-1)
  \;\Longrightarrow\;
  \frac{y-3}{2}=\ln(x-1)
  \;\Longrightarrow\;
  \mathrm{e}^{(y-3)/2}=x-1
  \;\Longrightarrow\;
  x=1+\mathrm{e}^{(y-3)/2} .
\]
Renaming the variable,
\[
  \boxed{\;g^{-1}(x)=1+\mathrm{e}^{(x-3)/2}\;}
\]

\wstep{Step 7. State the domain and range, which swap.}
$g$ has domain $(1,\infty)$ and range $\R$. Therefore $g^{-1}$ has domain
$\R$ and range $(1,\infty)$ --- visible in the formula, since
$\mathrm{e}^{(x-3)/2}>0$ always, so $g^{-1}>1$ always.
\wtrap{The domain of the inverse is the \emph{range} of the original. This is
not bookkeeping: it is what makes $g^{-1}$ a function at all. And check
$g^{-1}(g(x))=x$ on one value --- $g(1+\mathrm{e}^{0})=g(2)=2\ln1+3=3$ and
$g^{-1}(3)=1+\mathrm{e}^{0}=2$. \checkmark}
\end{worked}

\section{Practice Problems}

In problems 1--10, describe in words how the graph of $y=f(x)$ has to be
moved, scaled or reflected.

\begin{problems}
  \item $y=f(x-3)$
  \item $y=f(x)+2$
  \item $y=-f(x)$
  \item $y=f(-x)$
  \item $y=2f(x)$
  \item $y=f(2x)$
  \item $y=f\!\left(\dfrac{x}{2}\right)+1$
  \item $y=-3f(x-1)+4$
  \item $y=f(x+2)-5$
  \item $y=\tfrac12 f(-x)$
\end{problems}

In problems 11--18, write down the formula.

\begin{problems}[resume]
  \item $y=x^{2}$ shifted $3$ to the right and $1$ up.
  \item $y=\sqrt{x}$ reflected in the $x$-axis and shifted $4$ up.
  \item $y=\mathrm{e}^{x}$ compressed horizontally by the factor $2$.
  \item $y=\ln x$ shifted $5$ to the left.
  \item $y=1/x$ shifted $2$ to the right and $3$ down.
  \item $y=|x|$ shifted $1$ to the left and stretched vertically by $3$.
  \item $y=x^{3}$ reflected in the $y$-axis and shifted $2$ up.
  \item A downward parabola with vertex $(-1,5)$, twice as steep as
        $y=x^{2}$.
\end{problems}

In problems 19--24, name the base function and the transformations, and
give the features asked for.

\begin{problems}[resume]
  \item $y=2(x+1)^{2}-8$: vertex, roots, range.
  \item $y=-\sqrt{x-2}+3$: domain, starting point, range.
  \item $y=3\mathrm{e}^{x+1}$: asymptote, value at $x=-1$.
  \item $y=\ln(2x)$: give \emph{two} readings of this graph, one
        horizontal and one vertical.
  \item $y=\dfrac{1}{x-4}+2$: both asymptotes.
  \item $y=-|x+3|+5$: vertex, roots, direction of opening.
\end{problems}

\begin{problems}[resume]
  \item For $f(x)=x^{2}-2x-3$, describe the graph of $y=|f(x)|$: which
        part moves, and what happens at the roots?
  \item Sketch $y=|2x-6|$ by first writing it as a stretched, shifted V.
  \item For $f(x)=\sqrt{x}$, describe $y=f(|x|)$: give its domain and its
        symmetry.
  \item With $f(u)=u^{2}$ and $g(x)=x-5$: give $(f\circ g)(x)$ and
        $(g\circ f)(x)$, and say which transformation of $u^{2}$ each one
        is.
  \item With $f(u)=\sqrt{u}$ and $g(x)=x^{2}+9$: give $(f\circ g)(x)$ and
        its domain.
  \item Decompose $h(x)=\sqrt{3x+1}$ into $f\circ g$.
  \item Decompose $h(x)=\mathrm{e}^{x^{2}}$ into $f\circ g$.
  \item Decompose $h(x)=(\ln x)^{3}$ into $f\circ g$.
  \item Decompose $h(x)=\dfrac{1}{x^{2}+1}$ into $f\circ g$.
\end{problems}

In problems 34--39, find the inverse together with its domain and range.

\begin{problems}[resume]
  \item $f(x)=3x-7$
  \item $f(x)=\dfrac{x+2}{x-1}$
  \item $f(x)=x^{2}$ on $D=[0,\infty)$
  \item $f(x)=\ln(x-1)$
  \item $f(x)=\sqrt{x+4}$
  \item $f(x)=\dfrac{2x+1}{x-3}$
  \item Verify for $f(x)=3x-7$ that $f\bigl(f^{-1}(x)\bigr)=x$ and
        $f^{-1}\bigl(f(x)\bigr)=x$.
  \item The demand function is $q(p)=50-2p$. Give the inverse demand
        function $p(q)$, and state on which interval of quantities it is
        economically meaningful.
\end{problems}

\begin{problems}[resume]
  \item Determine $a$ so that the parabola with vertex $(2,-3)$ passes
        through $(0,1)$, and write the function.
  \item An exponential $y=A\,q^{x}$ passes through $(0,5)$ and $(2,45)$.
        Determine $A$ and $q$.
  \item Write a rational function with vertical asymptote $x=2$,
        horizontal asymptote $y=3$, and a root at the origin.
  \item Total cost is $C(q)=5q+100$. Write the new cost function and name
        the transformation when (a) the fixed cost rises by $40$;
        (b) a tax of $2$ per unit is levied.
\end{problems}

\section{Solutions}

\begin{solutions}
  \item Shifted $3$ units to the \emph{right} (inside, against the sign).
  \item Shifted $2$ units up.
  \item Reflected in the $x$-axis: every $y$-value changes sign.
  \item Reflected in the $y$-axis: left and right are swapped.
  \item Stretched vertically by the factor $2$; roots and $x$-intercepts
        stay where they are.
  \item Compressed horizontally by the factor $2$: what happened at $x=4$
        now happens at $x=2$. The $y$-intercept is unchanged.
  \item Stretched horizontally by the factor $2$ (the inner factor
        $\tfrac12$ is a compression by $\tfrac12$, i.e. a stretch), then
        moved $1$ up.
  \item Right $1$, stretched vertically by $3$, reflected in the $x$-axis,
        then up $4$ --- inside first, outside afterwards.
  \item Left $2$ and down $5$.
  \item Reflected in the $y$-axis and compressed vertically to half its
        height.
  \item $y=(x-3)^{2}+1$.
  \item $y=-\sqrt{x}+4$.
  \item $y=\mathrm{e}^{2x}$ (a factor $2$ \emph{inside} compresses).
  \item $y=\ln(x+5)$, with the vertical asymptote moved to $x=-5$.
  \item $y=\dfrac{1}{x-2}-3$, asymptotes $x=2$ and $y=-3$.
  \item $y=3|x+1|$.
  \item $y=(-x)^{3}+2=-x^{3}+2$.
  \item $y=-2(x+1)^{2}+5$.
  \item Base $x^{2}$: left $1$, stretched by $2$, down $8$. Vertex
        $(-1,-8)$; roots from $2(x+1)^{2}=8\iff(x+1)^{2}=4\iff x=1$ or
        $x=-3$; range $[-8,\infty)$.
  \item Base $\sqrt{x}$: right $2$, reflected in the $x$-axis, up $3$.
        $D=[2,\infty)$, the curve starts at $(2,3)$ and falls,
        $W=(-\infty,3]$.
  \item Base $\mathrm{e}^{x}$: left $1$, stretched by $3$. Asymptote
        $y=0$; $y(-1)=3\mathrm{e}^{0}=3$. (Equivalently
        $3\mathrm{e}\cdot\mathrm{e}^{x}$: a horizontal shift of an
        exponential is also a vertical stretch.)
  \item Horizontally: $\ln x$ compressed by the factor $2$. Vertically:
        since $\ln(2x)=\ln 2+\ln x$, the same graph is $\ln x$ shifted up
        by $\ln2\approx0.69$. Both readings are correct --- a peculiarity
        of the logarithm.
  \item Base $1/x$: right $4$, up $2$. Vertical asymptote $x=4$,
        horizontal asymptote $y=2$.
  \item Base $|x|$: left $3$, reflected (so it opens downwards), up $5$.
        Vertex $(-3,5)$; roots where $|x+3|=5$, i.e. $x=2$ and $x=-8$.
  \item $f(x)=(x-3)(x+1)$ is negative exactly on $(-1,3)$, and that arc is
        folded up above the axis; outside it the graph is unchanged. At the
        roots $x=-1$ and $x=3$ the curve now touches the axis and turns
        instead of crossing it, and the old minimum $f(1)=-4$ becomes a
        local maximum of height $4$.
  \item $|2x-6|=2|x-3|$: a V with vertex $(3,0)$, twice as steep as
        $|x|$, i.e. with slopes $-2$ and $+2$; $y$-intercept $6$.
  \item $y=\sqrt{|x|}$: $D=\R$ now, since $|x|\geq0$ always. The right half
        is unchanged and mirrored to the left, so the function is even,
        with range $[0,\infty)$ and a sharp point at the origin.
  \item $(f\circ g)(x)=(x-5)^{2}$: the parabola shifted $5$ to the right.
        $(g\circ f)(x)=x^{2}-5$: the parabola shifted $5$ down. Inner
        shifts move horizontally, outer shifts vertically.
  \item $(f\circ g)(x)=\sqrt{x^{2}+9}$, and since $x^{2}+9\geq9>0$ for
        every $x$, $D=\R$. (Its range is $[3,\infty)$.)
  \item $g(x)=3x+1$, $f(u)=\sqrt{u}$.
  \item $g(x)=x^{2}$, $f(u)=\mathrm{e}^{u}$.
  \item $g(x)=\ln x$, $f(u)=u^{3}$.
  \item $g(x)=x^{2}+1$, $f(u)=\dfrac{1}{u}$.
  \item $y=3x-7\Rightarrow x=\dfrac{y+7}{3}$, so
        $f^{-1}(x)=\dfrac{x+7}{3}$, with $D=W=\R$.
  \item $y=\dfrac{x+2}{x-1}\Rightarrow y(x-1)=x+2\Rightarrow x(y-1)=y+2
        \Rightarrow x=\dfrac{y+2}{y-1}$. Hence
        $f^{-1}(x)=\dfrac{x+2}{x-1}=f(x)$: the function is its own
        inverse, and its graph is symmetric about $y=x$. Domain and range
        are both $\R\setminus\{1\}$.
  \item $f^{-1}(x)=\sqrt{x}$, with $D(f^{-1})=[0,\infty)=W(f)$ and
        $W(f^{-1})=[0,\infty)=D(f)$. The restriction to $[0,\infty)$ was
        needed to make $f$ injective in the first place.
  \item $y=\ln(x-1)\Rightarrow \mathrm{e}^{y}=x-1\Rightarrow
        x=\mathrm{e}^{y}+1$: $f^{-1}(x)=\mathrm{e}^{x}+1$, with $D=\R$ and
        $W=(1,\infty)$ --- exactly the domain $(1,\infty)$ of $f$.
  \item $y=\sqrt{x+4}\Rightarrow y^{2}=x+4\Rightarrow x=y^{2}-4$, and only
        $y\geq0$ occurs: $f^{-1}(x)=x^{2}-4$ with $D=[0,\infty)$ and
        $W=[-4,\infty)$. Without the restriction $D=[0,\infty)$ the formula
        $x^{2}-4$ would not be the inverse.
  \item $y=\dfrac{2x+1}{x-3}\Rightarrow y(x-3)=2x+1\Rightarrow
        x(y-2)=3y+1\Rightarrow x=\dfrac{3y+1}{y-2}$. So
        $f^{-1}(x)=\dfrac{3x+1}{x-2}$, with $D=\R\setminus\{2\}$ (the
        range of $f$) and $W=\R\setminus\{3\}$ (the domain of $f$). The two
        asymptotes have swapped roles.
  \item $f\bigl(f^{-1}(x)\bigr)=3\cdot\dfrac{x+7}{3}-7=x+7-7=x$, and
        $f^{-1}\bigl(f(x)\bigr)=\dfrac{(3x-7)+7}{3}=\dfrac{3x}{3}=x$.
  \item $q=50-2p\Rightarrow 2p=50-q\Rightarrow p(q)=25-\dfrac{q}{2}$. It
        is meaningful for $q\in[0,50]$: at $q=50$ the price has fallen to
        $0$, and larger quantities would require a negative price.
  \item $y=a(x-2)^{2}-3$ with $1=a\cdot4-3$, so $a=1$ and
        $y=(x-2)^{2}-3$.
  \item $A=y(0)=5$. Then $45=5q^{2}\Rightarrow q^{2}=9\Rightarrow q=3$
        (positive base), so $y=5\cdot3^{x}$.
  \item $y=\dfrac{3x}{x-2}$: the denominator gives the pole $x=2$, equal
        degrees give the horizontal asymptote $y=3$, and the numerator
        vanishes at $x=0$, so the graph passes through the origin.
  \item (a) $C(q)=5q+140$: a pure vertical shift by $40$; the curve keeps
        its slope. (b) $C(q)=7q+100$: the tax adds $2q$, so the
        \emph{slope} changes from $5$ to $7$ while the intercept stays.
        A vertical shift and a change of slope are different operations,
        and only the first one is a translation.
\end{solutions}

\section{Additional Problems}

\begin{problems}
  \item $y=(x+4)^{2}-2$: base function, transformations, vertex, range.
  \item $y=-2\sqrt{x}+1$: transformations, domain, range.
  \item $y=\mathrm{e}^{-(x-1)}$: transformations, asymptote, value at
        $x=1$.
  \item $y=\ln(x+3)$: transformation, vertical asymptote, domain.
  \item $y=\dfrac{3}{x+2}-1$: transformations and both asymptotes.
  \item $y=\tfrac12|x-6|$: transformations and vertex.
  \item Write: $y=1/x$ reflected in the $x$-axis and shifted $3$ up.
  \item Write: $y=\mathrm{e}^{x}$ compressed horizontally by the factor
        $3$ and shifted $1$ down.
  \item Write: $y=\sqrt{x}$ stretched vertically by $5$ and shifted $1$ to
        the right.
  \item Write: a downward parabola with vertex $(3,-4)$, twice as steep as
        $y=x^{2}$.
  \item Describe the graph of $y=|x^{2}-1|$: which arc is folded, and what
        does the graph do at $x=\pm1$?
  \item Describe the graph of $y=|\ln x|$: domain, roots, and behaviour as
        $x\to0^{+}$.
  \item For $f(x)=x^{3}-3x$, give the formula for $f(|x|)$ and state its
        symmetry.
  \item With $f(u)=u^{2}$ and $g(x)=x+3$: give both composites and say why
        they differ.
  \item Decompose $h(x)=\sqrt{1-x^{2}}$ into $f\circ g$ and give its
        domain.
  \item Decompose $h(x)=\mathrm{e}^{\sqrt{x}}$ into $f\circ g$.
  \item Decompose $h(x)=(2x+5)^{10}$ into $f\circ g$.
  \item Inverse of $f(x)=5-2x$.
  \item Inverse of $f(x)=x^{3}+1$.
  \item Inverse of $f(x)=\dfrac{1}{x-2}$, with domain and range.
  \item Inverse of $f(x)=10^{x}$, with domain and range.
  \item Inverse of $f(x)=\ln(3x)$, with domain and range.
  \item Why does $f(x)=x^{2}$ have no inverse on $\R$, and which two
        restrictions repair it?
  \item A quantity is $200$ at $x=0$ and halves every $5$ units. Write it
        both as $A\,q^{x}$ and in the form $A\,\mathrm{e}^{-kx}$.
  \item Cost is $C(q)=3q+40$. Fixed cost rises to $60$ and the unit cost
        falls to $2.5$. Write the new function and name both changes as
        transformations of the graph.
\end{problems}

In problems 26--36, rewrite the formula first (single base, factored
inside, factored power), then give the chain of moves from the base
function, the maximal domain, the range and every asymptote.

\begin{problems}[resume]
  \item $y=3-2\cdot5^{-x}$
  \item $y=1+3\cdot\bigl(\tfrac25\bigr)^{x}$
  \item $y=\dfrac{12^{x}}{3^{x}}-2$
  \item $y=4-\dfrac{2^{2x}}{8^{x}}$
  \item $y=2\sqrt{6-2x}+1$
  \item $y=2-\sqrt{5-x}$
  \item $y=3\ln(2x+6)-1$
  \item $y=1-2\ln(4-x)$
  \item $y=\dfrac{2}{(x-1)^{2}}-3$
  \item $y=1-\dfrac{1}{(x+2)^{3}}$
  \item $y=\dfrac{4}{(2x-6)^{2}}+1$
\end{problems}

\section{Additional Solutions}

\begin{solutions}
  \item Base $x^{2}$, shifted $4$ to the left and $2$ down; vertex
        $(-4,-2)$; range $[-2,\infty)$.
  \item Base $\sqrt{x}$: stretched by $2$, reflected in the $x$-axis, then
        up $1$. $D=[0,\infty)$, and since $-2\sqrt{x}\in(-\infty,0]$,
        $W=(-\infty,1]$.
  \item $\mathrm{e}^{-(x-1)}$: reflect $\mathrm{e}^{x}$ in the $y$-axis,
        then shift $1$ to the right (the inner expression is
        $-(x-1)$). Asymptote $y=0$; $y(1)=\mathrm{e}^{0}=1$.
  \item $\ln x$ shifted $3$ to the left; the asymptote moves with it to
        $x=-3$; $D=(-3,\infty)$.
  \item Base $1/x$: stretched by $3$, left $2$, down $1$. Vertical
        asymptote $x=-2$, horizontal asymptote $y=-1$.
  \item Base $|x|$: right $6$ and compressed vertically to half the height;
        vertex $(6,0)$, slopes $\pm\tfrac12$.
  \item $y=-\dfrac{1}{x}+3$.
  \item $y=\mathrm{e}^{3x}-1$.
  \item $y=5\sqrt{x-1}$.
  \item $y=-2(x-3)^{2}-4$.
  \item The arc between the roots, where $x^{2}-1<0$, i.e. over
        $(-1,1)$, is folded up; it becomes a hump of height $1$ at $x=0$.
        At $x=\pm1$ the graph touches the axis and turns, forming two
        corners.
  \item $D=(0,\infty)$; the single root is $x=1$; on $(0,1)$ the logarithm
        is negative and is folded up, so the graph falls from $+\infty$ to
        $0$ there and rises again afterwards, with a corner at $(1,0)$. As
        $x\to0^{+}$ it goes to $+\infty$.
  \item $f(|x|)=|x|^{3}-3|x|$. It is even: the right half of the original
        cubic is kept and mirrored into the left half-plane, so the old
        left branch is discarded.
  \item $(f\circ g)(x)=(x+3)^{2}$, the parabola shifted $3$ to the left;
        $(g\circ f)(x)=x^{2}+3$, the parabola shifted $3$ up. Composition
        is not commutative: an inner shift moves the graph horizontally, an
        outer one vertically.
  \item $g(x)=1-x^{2}$, $f(u)=\sqrt{u}$; $D=[-1,1]$ from $1-x^{2}\geq0$.
  \item $g(x)=\sqrt{x}$, $f(u)=\mathrm{e}^{u}$; $D=[0,\infty)$.
  \item $g(x)=2x+5$, $f(u)=u^{10}$.
  \item $y=5-2x\Rightarrow x=\dfrac{5-y}{2}$, so
        $f^{-1}(x)=\dfrac{5-x}{2}$.
  \item $y=x^{3}+1\Rightarrow x=\sqrt[3]{y-1}$, so
        $f^{-1}(x)=\sqrt[3]{x-1}$, with $D=W=\R$.
  \item $y=\dfrac{1}{x-2}\Rightarrow x-2=\dfrac1y\Rightarrow
        x=2+\dfrac1y$: $f^{-1}(x)=2+\dfrac{1}{x}$, with
        $D=\R\setminus\{0\}$ and $W=\R\setminus\{2\}$.
  \item $f^{-1}(x)=\log_{10}x$, with $D=(0,\infty)$ and $W=\R$.
  \item $y=\ln(3x)\Rightarrow \mathrm{e}^{y}=3x\Rightarrow
        x=\dfrac{\mathrm{e}^{y}}{3}$: $f^{-1}(x)=\dfrac{\mathrm{e}^{x}}{3}$,
        with $D=\R$ and $W=(0,\infty)$.
  \item It is not injective: $f(-2)=f(2)=4$, so a horizontal line meets the
        graph twice and the inverse assignment would not be unique.
        Restricting to $[0,\infty)$ gives the inverse $\sqrt{x}$;
        restricting to $(-\infty,0]$ gives $-\sqrt{x}$.
  \item Halving every $5$ units means $q=2^{-1/5}$, so
        $y=200\cdot2^{-x/5}$. Since $2^{-x/5}
        =\mathrm{e}^{-(\ln2/5)x}$, the same function is
        $y=200\,\mathrm{e}^{-kx}$ with
        $k=\dfrac{\ln2}{5}\approx0.1386$.
  \item $C(q)=2.5q+60$. Raising the fixed cost from $40$ to $60$ shifts the
        line $20$ units up; lowering the unit cost from $3$ to $2.5$
        flattens it, a change of slope rather than a translation. Only the
        first is a shift of the graph.
  \item Rewrite $5^{-x}=\bigl(\tfrac15\bigr)^{x}$. From $y=5^{x}$:
        reflect in the $y$-axis (the graph now falls), stretch vertically by
        $2$, reflect in the $x$-axis (it now rises again), shift up $3$.
        $D=\R$. Since $2\cdot5^{-x}$ takes every value in $(0,\infty)$,
        $-2\cdot5^{-x}$ takes every value in $(-\infty,0)$ and
        $W=(-\infty,3)$. Horizontal asymptote $y=3$, approached from below
        as $x\to\infty$; as $x\to-\infty$ the graph falls without bound.
        Check: $y(0)=3-2=1$. Root: $5^{-x}=\tfrac32$, so
        $x=-\dfrac{\ln1.5}{\ln5}\approx-0.25$.
  \item The base $\tfrac25$ lies between $0$ and $1$, so $(\tfrac25)^{x}$
        is a \emph{falling} exponential through $(0,1)$. Stretch vertically
        by $3$, shift up $1$. $D=\R$, $W=(1,\infty)$, horizontal asymptote
        $y=1$, approached from above as $x\to\infty$. $y(0)=1+3=4$.
        (Equivalently $3\cdot(\tfrac52)^{-x}+1$: the rising exponential
        $(\tfrac52)^{x}$ reflected in the $y$-axis.)
  \item Rewrite $\dfrac{12^{x}}{3^{x}}=\Bigl(\dfrac{12}{3}\Bigr)^{x}=4^{x}$,
        so $y=4^{x}-2$: the rising exponential $4^{x}$ shifted $2$ down.
        $D=\R$, $W=(-2,\infty)$, horizontal asymptote $y=-2$ as
        $x\to-\infty$. $y(0)=-1$; root $4^{x}=2$, i.e. $x=\tfrac12$.
  \item Rewrite $2^{2x}=4^{x}$, so $\dfrac{2^{2x}}{8^{x}}
        =\Bigl(\dfrac48\Bigr)^{x}=\bigl(\tfrac12\bigr)^{x}=2^{-x}$ and
        $y=4-2^{-x}$. From $2^{x}$: reflect in the $y$-axis, reflect in the
        $x$-axis, shift up $4$. $D=\R$, $W=(-\infty,4)$, horizontal
        asymptote $y=4$ as $x\to\infty$. $y(0)=3$; root $2^{-x}=4=2^{2}$,
        so $x=-2$.
  \item Factor out $-2$: $6-2x=-2(x-3)$, so $y=2\sqrt{-2(x-3)}+1$. From
        $\sqrt{x}$: reflect in the $y$-axis and compress horizontally by
        $2$ (giving $\sqrt{-2x}$), shift right $3$, stretch vertically by
        $2$, shift up $1$. Domain from $6-2x\ge0$: $D=(-\infty,3]$. The
        curve starts at $(3,1)$ and rises to the left without bound, so
        $W=[1,\infty)$; no asymptote. Check: $y(1)=2\sqrt4+1=5$.
  \item Factor out $-1$: $y=-\sqrt{-(x-5)}+2$. From $\sqrt{x}$: reflect in
        the $y$-axis, shift right $5$, reflect in the $x$-axis, shift up
        $2$. $D=(-\infty,5]$, starting point $(5,2)$, from which the curve
        falls to the left: $W=(-\infty,2]$. No asymptote. Root:
        $\sqrt{5-x}=2$, so $x=1$.
  \item Factor the inside: $2x+6=2(x+3)$, so $y=3\ln\bigl(2(x+3)\bigr)-1$.
        From $\ln x$: compress horizontally by $2$, shift left $3$, stretch
        vertically by $3$, shift down $1$. $D=(-3,\infty)$ from $2x+6>0$;
        vertical asymptote $x=-3$; $W=\R$. Second reading:
        $y=3\ln(x+3)+(3\ln2-1)$, i.e. $\ln(x+3)$ stretched by $3$ and
        shifted up by $3\ln2-1\approx1.08$. Root: $\ln(2x+6)=\tfrac13$,
        $x=\tfrac12\bigl(\mathrm{e}^{1/3}-6\bigr)\approx-2.30$.
  \item Factor out $-1$: $y=-2\ln\bigl(-(x-4)\bigr)+1$. From $\ln x$:
        reflect in the $y$-axis, shift right $4$, stretch vertically by
        $2$, reflect in the $x$-axis, shift up $1$. $D=(-\infty,4)$ from
        $4-x>0$; vertical asymptote $x=4$. As $x\to4^{-}$,
        $\ln(4-x)\to-\infty$, so $y\to+\infty$; the two reflections together
        leave an \emph{increasing} curve. $W=\R$. Root: $\ln(4-x)=\tfrac12$,
        $x=4-\sqrt{\mathrm{e}}\approx2.35$.
  \item Base $1/x^{2}$ (even power): stretch vertically by $2$, shift
        right $1$, down $3$. $D=\R\setminus\{1\}$, vertical asymptote $x=1$,
        horizontal asymptote $y=-3$. Since $\dfrac{2}{(x-1)^{2}}>0$ on both
        sides, both branches lie above $y=-3$: $W=(-3,\infty)$. Roots:
        $(x-1)^{2}=\tfrac23$, $x=1\pm\sqrt{2/3}\approx1\pm0.82$.
  \item Base $1/x^{3}$ (odd power): reflect in the $x$-axis, shift left
        $2$, up $1$. $D=\R\setminus\{-2\}$, asymptotes $x=-2$ and $y=1$.
        For $x>-2$ the cube $(x+2)^{3}$ is positive and $y<1$; for $x<-2$
        it is negative and $y>1$: the branches lie on opposite sides, and
        $W=\R\setminus\{1\}$. Root: $(x+2)^{3}=1$, $x=-1$; $y(0)=\tfrac78$.
  \item Factor the power: $(2x-6)^{2}=\bigl(2(x-3)\bigr)^{2}=4(x-3)^{2}$,
        so $y=\dfrac{1}{(x-3)^{2}}+1$. Base $1/x^{2}$: shift right $3$, up
        $1$. $D=\R\setminus\{3\}$, asymptotes $x=3$ and $y=1$, both branches
        above $y=1$: $W=(1,\infty)$. No roots.
\end{solutions}
